\documentclass[11pt]{article}
\usepackage{amsmath,amsthm,amssymb}
\usepackage[margin=1in]{geometry}

\newtheorem{theorem}{Theorem}
\newtheorem{proposition}{Proposition}
\newtheorem{lemma}{Lemma}
\newtheorem{corollary}{Corollary}
\newtheorem{conjecture}{Conjecture}
\theoremstyle{definition}
\newtheorem{remark}{Remark}

\newcommand{\Z}{\mathbb{Z}}
\newcommand{\C}{\mathbb{C}}
\newcommand{\R}{\mathbb{R}}
\newcommand{\Zi}{\Z[i]}
\renewcommand{\Re}{\operatorname{Re}}
\renewcommand{\Im}{\operatorname{Im}}

\title{The two-row $d=4$ fibre-vanishing law reduces to a single phase positivity:\\
$G_{(a,b)}(i)\neq0$ except $(2,2)$, via $\Im(r_b\overline{r_{b-1}})>0$}
\author{Clio}
\date{2026-06-04 (prove session)}

\begin{document}
\maketitle

\begin{abstract}
Let $f(u)=u^2+(1+i)u+1\in\Zi[u]$, $r_l=[u^l]f(u)^m$, and for a two-row partition
$(a,b)\vdash 2m$ ($a\ge b\ge1$) let $G_{(a,b)}(i)=r_b-r_{b-1}$.  The two-row case of
the $d=4$ fibre-vanishing trichotomy asserts $G_{(a,b)}(i)\neq0$ for all two-row
$(a,b)$ \emph{except} $(2,2)$.  We replace the prior route (real monotonicity of
$N_l=|r_l|^2$, which stalled at a single endpoint inequality $D_m>0$) by a
\emph{uniform} reduction covering all two-row shapes at once.  Writing
$C_k=r_k\overline{r_{k-1}}=P_k+iQ_k$, we show
\[
   Q_b\neq0\ \Longrightarrow\ G_{(a,b)}(i)\neq0 .
\]
We prove $P_k>0$ for all $1\le k\le m$ unconditionally, derive the exact
Casoratian recurrence $(k+1)Q_{k+1}=(m-k)N_k-(2m-k+1)Q_k$, and reduce the entire
two-row law to the \emph{phase-positivity statement}
\[
   (\mathrm{PP})\qquad Q_k>0\quad\text{for }1\le k\le m,\ \ m\ge4,
\]
together with three finite checks ($m\le3$).  We isolate $(\mathrm{PP})$ to the
single inequality family $(B_k)\colon \beta_kQ_k<\alpha_kN_k$, exhibit the exact
two-dimensional dynamical system governing it, and identify the mechanism: in the
distance-from-top variable $j=m-k$ the binding ratio
$x_k=\beta_kQ_k/(\alpha_kN_k)$ obeys a limiting recursion
$y_{j-1}=\tfrac{j}{j-1}(1-y_j)$ whose self-propagating invariant
$\tfrac12<y_j<\tfrac{j+1}{2j}$ forces $x_k\to\tfrac34<1$.  All statements are
verified exactly/at $50$-digit precision for $m\le419$.  The one remaining gap is
precisely a uniform-in-$m$ control of the corrections to this limiting recursion.
\end{abstract}

\section{Setup and the new reduction}

Fix $i=\sqrt{-1}$ and an integer $m\ge1$.  Put
\[
  f(u)=u^2+(1+i)u+1\in\Zi[u],\qquad r_l:=[u^l]f(u)^m\quad(0\le l\le2m).
\]
Since $u^2f(1/u)=f(u)$ the coefficient sequence is palindromic, $r_l=r_{2m-l}$.
For a two-row partition $(a,b)\vdash2m$ with $a\ge b\ge1$ (so $1\le b\le m$), the
$\zeta_4=i$ fibre value is
\[
  G_{(a,b)}(i)=\langle s_{(a,b)},\psi^m\rangle\Big|_{\psi=h_2+ie_2}=r_b-r_{b-1},
\]
the second equality being the established two-row reduction.  The target of this
session is:

\begin{theorem}[two-row $d=4$ fibre vanishing]\label{thm:main}
For every $m\ge1$ and every two-row $(a,b)\vdash2m$ one has $G_{(a,b)}(i)\neq0$,
with the single exception $(a,b)=(2,2)$ (where $r_2=r_1$, $G=0$).
\end{theorem}

Throughout set
\[
  N_l:=|r_l|^2=r_l\overline{r_l}\in\Z_{\ge0},\qquad
  C_k:=r_k\overline{r_{k-1}}=P_k+iQ_k\ \ (k\ge1),
\]
with $P_k=\Re C_k$, $Q_k=\Im C_k$.  We also use $\alpha_k:=m-k$,
$\beta_k:=2m-k+1$ (so $\alpha_k\ge0$, $\beta_k>0$ for $0\le k\le m$).

The key new observation is that the truth of Theorem~\ref{thm:main} for the shape
$(a,b)$ follows from a single \emph{sign} condition on $Q_b$.

\begin{lemma}[phase criterion]\label{lem:phase}
If $Q_b\neq0$ then $G_{(a,b)}(i)=r_b-r_{b-1}\neq0$.
\end{lemma}
\begin{proof}
Contrapositive.  If $r_b=r_{b-1}$ then
$C_b=r_b\overline{r_{b-1}}=r_{b-1}\overline{r_{b-1}}=|r_{b-1}|^2\in\R$, hence
$Q_b=\Im C_b=0$.
\end{proof}

Thus Theorem~\ref{thm:main} is implied by the positivity of the imaginary parts
$Q_b$.  We record what actually happens.

\begin{proposition}[reduction to phase positivity]\label{prop:reduction}
Suppose the following \emph{phase-positivity statement} holds:
\[
  (\mathrm{PP})\qquad Q_k>0\quad\text{for all } 1\le k\le m \text{ and all } m\ge4 .
\]
Then Theorem~\ref{thm:main} holds.
\end{proposition}
\begin{proof}
For $m\ge4$, $(\mathrm{PP})$ gives $Q_b>0$ for every $1\le b\le m$, so
Lemma~\ref{lem:phase} yields $G_{(a,b)}(i)\neq0$ for all two-row shapes.  For
$m\le3$ there are only finitely many two-row shapes; direct evaluation in $\Zi$
gives, with $z=-1/(3+i)$ and the closed form of \S\ref{sec:saddle},
\[
\begin{array}{llll}
 m=1: & G_{(1,1)}=i\neq0; \\[2pt]
 m=2: & G_{(3,1)}=1+2i\neq0, & G_{(2,2)}=0 \ (\text{the exception}); \\[2pt]
 m=3: & G_{(5,1)}=2+3i,\ & G_{(4,2)}=3i,\ & G_{(3,3)}=1+2i\ (\text{all}\neq0).
\end{array}
\]
(For $m=3$ one has $Q_3=0$, so Lemma~\ref{lem:phase} does \emph{not} apply at
$(3,3)$; the value $G_{(3,3)}=(3+i)^3\cdot\Sigma_0$ with $\Sigma_0=7/100+i/100$ is
nonzero by direct computation.)  Hence the only vanishing fibre value is at
$(2,2)$.
\end{proof}

\begin{remark}
The point of the reduction is that it treats \emph{all} two-row shapes uniformly
through one positivity, in contrast to the earlier real-monotonicity route
(\,$N_0<N_1<\cdots<N_m$ with the lone tie $N_1=N_2$ at $m=2$\,), which singled out
the balanced shape $(m,m)$ as a hard endpoint inequality $D_m=N_m-N_{m-1}>0$ left
unproven for general $m$.  The sign $Q_b\neq0$ is logically weaker than
$|r_b|\neq|r_{b-1}|$, yet still suffices, and turns out to be the natural quantity
for the recurrence.
\end{remark}

\section{An exact saddle expansion}\label{sec:saddle}

The constant $3+i$ that appears repeatedly originates in an exact algebraic
identity.  Since $f(1)=3+i$ and $f'(u)=2u+(1+i)$ gives $f'(1)=3+i=f(1)$, one has
the exact factor-free expansion about $u=1$:
\begin{equation}\label{eq:saddle}
  f(u)=(u-1)^2+(3+i)\,u .
\end{equation}
Expanding $f(u)^m=(3+i)^mu^m\big(1+\tfrac{(u-1)^2}{(3+i)u}\big)^m$ and extracting
coefficients gives, for $0\le s=m-b$,
\[
  r_{m-s}=(3+i)^m(-1)^s\sum_{j=0}^m\binom mj\binom{2j}{j-s}\frac{(-1)^j}{(3+i)^j},
\]
and hence, using $\binom{2j}{j-s}+\binom{2j}{j-s-1}=\binom{2j+1}{j-s}$,
\begin{equation}\label{eq:Greform}
  G_{(a,b)}(i)=r_b-r_{b-1}=(3+i)^m(-1)^{m-b}\,\Sigma_{m-b},\qquad
  \Sigma_s:=\sum_{j=0}^m\binom mj\binom{2j+1}{j-s}z^j,\ \ z=\tfrac{-1}{3+i}.
\end{equation}
Thus Theorem~\ref{thm:main} is also equivalent to $\Sigma_s\neq0$ for all
$0\le s\le m-1$, $m\ge1$, except $(m,s)=(2,0)$.  Identity \eqref{eq:Greform} was
verified symbolically for all $m\le8$.  (We do not use \eqref{eq:Greform} below
beyond the $m\le3$ checks; it is recorded because it makes the exceptional shape
$(2,2)$ and the constant $3+i$ transparent, and because it shows the
$(1+i)$-adic route is hopeless: $v_{1+i}(G_{(a,b)})$ is finite but unbounded in
$m$, e.g.\ $11$ at $(10,10)$.)

\section{The Casoratian recurrence and $P_k>0$}

The recurrence for $r_l$ (correct and verified in $\Zi$),
\begin{equation}\label{eq:rrec}
  (k+1)r_{k+1}=(1+i)\alpha_k\,r_k+\beta_k\,r_{k-1},\qquad r_{-1}=0,\ r_0=1,
\end{equation}
descends to an exact recurrence for $C_k=r_k\overline{r_{k-1}}$.

\begin{proposition}[bilinear recurrence]\label{prop:Crec}
For $1\le k\le m-1$,
\begin{equation}\label{eq:Crec}
  (k+1)\,C_{k+1}=(1+i)\alpha_k N_k+\beta_k\,\overline{C_k},
\end{equation}
equivalently, taking real and imaginary parts,
\begin{equation}\label{eq:PQrec}
  (k+1)P_{k+1}=\alpha_k N_k+\beta_k P_k,\qquad
  (k+1)Q_{k+1}=\alpha_k N_k-\beta_k Q_k .
\end{equation}
\end{proposition}
\begin{proof}
Multiply \eqref{eq:rrec} by $\overline{r_k}$:
$(k+1)r_{k+1}\overline{r_k}=(1+i)\alpha_k|r_k|^2+\beta_k r_{k-1}\overline{r_k}$.
Now $r_{k+1}\overline{r_k}=C_{k+1}$, $|r_k|^2=N_k$, and
$r_{k-1}\overline{r_k}=\overline{r_k\overline{r_{k-1}}}=\overline{C_k}$, giving
\eqref{eq:Crec}.  Separating real/imaginary parts, and using
$\Re\overline{C_k}=P_k$, $\Im\overline{C_k}=-Q_k$, yields \eqref{eq:PQrec}.
\end{proof}

Equation \eqref{eq:Crec} is the Casoratian (discrete Wronskian) of the pair
$(r_l,\overline{r_l})$: indeed $w_k:=r_k\overline{r_{k-1}}-\overline{r_k}r_{k-1}
=2iQ_k$ satisfies $(k+1)w_{k+1}=2i\alpha_kN_k-\beta_k w_k$, which is
\eqref{eq:PQrec} for $Q$.  All identities verified exactly in $\Zi$ for
$1\le m\le40$.

\begin{lemma}[real-part positivity]\label{lem:Ppos}
For all $1\le k\le m$ one has $P_k>0$; in particular $r_k\neq0$ and $N_k>0$.
\end{lemma}
\begin{proof}
Induction on $k$.  Base: $r_0=1$, $r_1=m(1+i)$, so $C_1=m(1+i)$ and $P_1=m>0$.
Inductive step: assume $P_k>0$.  Then $C_k=r_k\overline{r_{k-1}}\neq0$, so
$r_k\neq0$ and $r_{k-1}\neq0$; in particular $N_k=|r_k|^2>0$.  By the first
identity of \eqref{eq:PQrec}, $(k+1)P_{k+1}=\alpha_kN_k+\beta_kP_k$.  For
$k\le m-1$ we have $\alpha_k\ge0$, $N_k>0$, $\beta_k>0$, $P_k>0$, so each summand
is $\ge0$ and the second is $>0$; hence $P_{k+1}>0$.  (At $k=m$ the claim
$P_m>0$ has already been produced at the previous step.)
\end{proof}

\section{Reduction of phase positivity to a single inequality}

By the second identity of \eqref{eq:PQrec}, $Q_{k+1}>0$ is \emph{equivalent} to
\begin{equation}\label{eq:Bk}
  (B_k)\qquad \beta_k Q_k<\alpha_k N_k .
\end{equation}
Since $Q_1=m>0$ always, statement $(\mathrm{PP})$ ($Q_k>0$ for $1\le k\le m$) is
\emph{equivalent} to the family $(B_k)$ holding for all $1\le k\le m-1$.  Writing
\[
  x_k:=\frac{\beta_kQ_k}{\alpha_kN_k}\qquad(1\le k\le m-1,\ \alpha_k\ge1),
\]
the family $(B_k)$ is exactly $x_k<1$ for all $1\le k\le m-1$.

\begin{lemma}[base cases of $(B_k)$]\label{lem:base}
$x_1=\dfrac{1}{m-1}$.  Hence $x_1<1$ (so $Q_2>0$) for all $m\ge3$, while
$x_1=1$ (so $Q_2=0$) exactly at $m=2$.  More precisely $Q_2=m^2(m-2)$.
\end{lemma}
\begin{proof}
$r_1=m(1+i)$, $N_1=2m^2$, $Q_1=\Im\big(m(1+i)\big)=m$, $\alpha_1=m-1$,
$\beta_1=2m$, so $x_1=\frac{2m\cdot m}{(m-1)2m^2}=\frac1{m-1}$.  Then
$2Q_2=\alpha_1N_1-\beta_1Q_1=(m-1)2m^2-2m\cdot m=2m^2(m-2)$.
\end{proof}

So the whole of Theorem~\ref{thm:main} reduces to a single uniform statement:

\begin{conjecture}[bound $B$; the remaining gap]\label{conj:B}
For all $m\ge4$ and all $1\le k\le m-1$, $\ x_k=\dfrac{\beta_kQ_k}{\alpha_kN_k}<1$.
\end{conjecture}

Numerically (50-digit, $m\le419$): $x_k<1$ holds for all $m\ge3$ and
$1\le k\le m-1$; the supremum is approached at the top index $k=m-1$, where
$x_{m-1}\uparrow\frac34$ as $m\to\infty$.  The only values $x_k\ge1$ that occur are
the two genuine equalities $x_1=1$ at $m=2$ (the shape $(2,2)$) and
$x_2=1$ at $m=3$ (which records the harmless zero $Q_3=0$ at the shape $(3,3)$,
where $G\neq0$ nonetheless).  For $m\ge4$ one has $x_k\le 0.78$ throughout, the
maximum $0.7798$ being attained at $(m,k)=(6,5)$.

\section{The exact $2$-D system and the mechanism behind Conjecture~\ref{conj:B}}

We exhibit the dynamical system controlling $x_k$ and the structural reason it
stays below $1$.  Besides \eqref{eq:PQrec} we use the exact Gram relation (a
consequence of \eqref{eq:rrec} times its conjugate; verified to $m=40$):
\begin{equation}\label{eq:Nrec}
  (k+1)^2N_{k+1}=2\alpha_k^2N_k+\beta_k^2N_{k-1}+2\alpha_k\beta_k S_k,\qquad
  S_k:=P_k-Q_k,
\end{equation}
and the Cauchy--Schwarz \emph{equality} $P_k^2+Q_k^2=N_kN_{k-1}$ (valid because
$C_k$ is a genuine product).  Put $\mu_k:=N_k/N_{k-1}>0$ and
$\eta_k:=Q_k/N_k=\frac{\alpha_k}{\beta_k}x_k$.  Since $P_k>0$
(Lemma~\ref{lem:Ppos}), $P_k=N_k\sqrt{1/\mu_k-\eta_k^2}$ and
\eqref{eq:PQrec}--\eqref{eq:Nrec} close into the exact system
\begin{equation}\label{eq:system}
  x_{k+1}=\frac{\beta_{k+1}\alpha_k}{\alpha_{k+1}(k+1)}\cdot\frac{1-x_k}{\mu_{k+1}},
  \qquad
  (k+1)^2\mu_{k+1}=2\alpha_k^2+\frac{\beta_k^2}{\mu_k}
     +2\alpha_k\beta_k\!\left(\sqrt{\tfrac1{\mu_k}-\eta_k^2}-\eta_k\right).
\end{equation}
The first equation shows the structural reason $x$ resists exceeding $1$: it is of
the form $x_{k+1}=a_k(1-x_k)$ with $a_k>0$, whose only fixed point
$x^\ast=a_k/(1+a_k)$ is automatically $<1$.  The difficulty is purely that $a_k>1$
near the top (where $\alpha_{k+1}=m-k-1$ is small), so the bound $x_{k+1}<1$ needs
the factor $(1-x_k)$ to be correspondingly small, i.e.\ a matching \emph{lower}
bound on $x_k$.

\paragraph{Limiting profile.}  Let $j=m-k$ be the distance from the top and
$y_j:=x_{m-j}$.  Using $\mu_k\to1$, $\alpha_k=j$, $\beta_k\sim m$ as $m\to\infty$
with $j$ fixed, the first equation of \eqref{eq:system} degenerates to the clean
limiting recursion
\begin{equation}\label{eq:limit}
  y_{j-1}=\frac{j}{j-1}\,(1-y_j),
\end{equation}
which is matched by the data to $4$ decimals near the top.  Recursion
\eqref{eq:limit} possesses an exact self-propagating invariant:

\begin{lemma}[invariant for \eqref{eq:limit}]\label{lem:invariant}
If $\tfrac12<y_j<\tfrac{j+1}{2j}$ then $\tfrac12<y_{j-1}<\tfrac{j}{2(j-1)}
=\tfrac{(j-1)+1}{2(j-1)}$.  Consequently the band $\tfrac12<y_j<\tfrac{j+1}{2j}$
propagates from large $j$ down to $j=1$, where it gives $y_1<1$ and indeed
$y_1\to\tfrac34$ (with $y_\infty=\tfrac12$).
\end{lemma}
\begin{proof}
Upper: $y_j>\tfrac12\Rightarrow1-y_j<\tfrac12\Rightarrow
y_{j-1}=\tfrac{j}{j-1}(1-y_j)<\tfrac{j}{2(j-1)}$.  Lower:
$y_j<\tfrac{j+1}{2j}\Rightarrow1-y_j>\tfrac{j-1}{2j}\Rightarrow
y_{j-1}>\tfrac{j}{j-1}\cdot\tfrac{j-1}{2j}=\tfrac12$.  The upper bound at index
$j-1$ is $\tfrac{(j-1)+1}{2(j-1)}=\tfrac{j}{2(j-1)}$, which is exactly what was
obtained; the band is therefore invariant.  At $j=1$ the upper bound is
$\tfrac{2}{2}=1$, and strictness $y_1<1$ follows from $y_2>\tfrac12$.
\end{proof}

\begin{remark}[what remains]\label{rem:gap}
Lemma~\ref{lem:invariant} proves Conjecture~\ref{conj:B} for the limiting
recursion \eqref{eq:limit}.  The exact $x_k$ does \emph{not} satisfy the band
$\tfrac12<x_k$ globally (in the bulk $x_k$ dips below $\tfrac12$; e.g.\ at $m=40$
already $x_{33}=0.4927$); the band is only an attractor near the top.  The single
remaining step to a complete proof is therefore a \emph{uniform-in-$m$ control of
the corrections} in \eqref{eq:system} to the limiting recursion
\eqref{eq:limit}, valid on the finite top window where $a_k>1$ (a window of
bounded size in $j$), together with an elementary bulk estimate $x_k\le\tfrac12$
for $j=m-k$ large where $a_k<1$ makes $x_{k+1}=a_k(1-x_k)<a_k<1$ immediate once
$\mu_{k+1}\ge\frac{\beta_{k+1}\alpha_k}{\alpha_{k+1}(k+1)}$.  Both pieces are
inequalities in $(\mu_k,x_k)$ only; the obstruction to writing them down cleanly
is the square-root coupling through $\mu$ in \eqref{eq:system}.
\end{remark}

\section{Status}
\begin{itemize}
\item \textbf{Proved unconditionally:} the phase criterion (Lemma~\ref{lem:phase}),
the saddle reformulation \eqref{eq:saddle}--\eqref{eq:Greform}, the bilinear
recurrence \eqref{eq:Crec}--\eqref{eq:PQrec}, real-part positivity $P_k>0$
(Lemma~\ref{lem:Ppos}), the equivalence $(\mathrm{PP})\Leftrightarrow(B_k)\
\forall k$, the base cases (Lemma~\ref{lem:base}), the finite cases $m\le3$, and
the invariant for the limiting recursion (Lemma~\ref{lem:invariant}).
\item \textbf{Reduces Theorem~\ref{thm:main} to:} Conjecture~\ref{conj:B}
($x_k<1$ for $m\ge4$), equivalently phase positivity $(\mathrm{PP})$.
\item \textbf{Verified:} all statements exactly / at $50$ digits for $m\le419$;
$\sup_k x_k<1$ with the only equalities at $(m,k)\in\{(2,1),(3,2)\}$.
\item \textbf{Remaining gap (precise):} the uniform correction control of
Remark~\ref{rem:gap}.
\end{itemize}

\end{document}
